\section{Conclusion and Future Work}\label{sec:conclusion}

This paper presented a Kubernetes-native framework for managing embedded WebAssembly workloads across the Cloud\textendash{}Fog\textendash{}Edge continuum without requiring constrained devices to become Kubernetes nodes. The contribution integrates CRD-based intent, gateway-mediated TLS/CBOR communication, Zephyr firmware, WAMR execution, and evidence-driven reconciliation into a single verifiable operational path suitable for sustainable digital research infrastructures.

% Removed: The experimental evaluation on a K3s test bed with emulated embedded targets confirms that the architecture supports a complete enrollment-to-deployment workflow under repeatable conditions, with 100/100 transactional success and {\color{red}approximately 1.29~s mean end-to-end orchestration latency}. Measured southbound control traffic, firmware footprint, and fog-layer resource use provide quantitative anchors for subsequent energy-aware placement and green-experimentation studies.
{\color{red}The experimental evaluation on a K3s testbed with emulated embedded targets confirms that the architecture supports a complete enrollment-to-deployment workflow under repeatable conditions, with 100/100 transactional success and approximately 1.29~s mean end-to-end orchestration latency. The measured application-layer traffic, firmware footprint, and fog-layer resource use provide quantitative anchors for subsequent energy-aware placement, carbon-aware annotation, and green-experimentation studies, while direct energy optimization remains future work.}
\revD{The computational-energy pass of Section~\ref{sec:energy-instrumentation} should be
read with the same scope: it instruments a Wasmtime runtime at the fog layer, not the
Zephyr/WAMR firmware on the embedded target, and its Joule and carbon figures are proxies
derived from measured CPU-time with an external power coefficient. No energy measurement
on the constrained endpoint itself is claimed.}

Future work will extend the platform toward multi-gateway deployments, fault injection, physical hardware, wattmeter- or RAPL-validated energy measurements{\color{blue}, carbon-intensity annotation beyond the illustrative, derived-estimate pass reported in Section~\ref{sec:energy-instrumentation}}, and larger device populations.
